\section{Conclusion}

The paper asked why relativisation helps for some KPIs and harms others, how a less efficient relative feature can still inform, and whether the same geometry appears beyond sport.

The Paired Efficiency Factor $\eta=(1+\kappa)/(1+\kappa-2\sqrt{\kappa}\,\rho)$ extends Fisher's paired efficiency to unequal variances. The formula is distribution-free. Under bivariate normality, $\eta$ and information content are related but distinct, which is why a relative feature can improve prediction when $\eta<1$.

The sports inventory occupies all four $(\kappa,\rho)$ quadrants and is Q3-heavy. Four exemplars at comparable signal strength confirm both help and harm (\cref{tab:exemplars}). The same geometry appears, with different mean $\hat\eta$, in healthcare, clinical genomics, finance, and manufacturing (\cref{tab:validation}). Read with $\delta/\sigma_{\mathrm{A}}$, the taxonomy is a prior for when relative features should be used and when they should not.
